\documentclass[11pt,a4paper]{article}
\usepackage[T2A]{fontenc}
\usepackage[utf8]{inputenc}
\usepackage[russian]{babel}
\usepackage{geometry}
\usepackage{amsmath,amssymb,mathtools}
\usepackage{array,longtable,booktabs}
\usepackage{xcolor}
\usepackage{tikz}
\usepackage{hyperref}
\geometry{margin=20mm}
\hypersetup{colorlinks=true,linkcolor=blue!50!black}
\setlength{\parindent}{0pt}
\setlength{\parskip}{5pt}
\newcommand{\Part}[1]{\par\medskip\noindent\textbf{#1}\par\smallskip}
\newcommand{\Task}[2]{\clearpage\phantomsection\label{task:#1}\section*{#2}}
\newcommand{\Err}[1]{\textcolor{red}{Ошибка:} #1}
\newcommand{\Idea}[1]{\textcolor{blue!60!black}{Ключевая идея:} #1}
\newcommand{\Result}[1]{\textcolor{teal!60!black}{Итог:} #1}
\begin{document}
\begin{titlepage}
\centering\vspace*{0.30\textheight}
{\huge\bfseries Ревью домашней работы\par}
\vspace{24mm}{\large Автор: Лёвин Артём Александрович\par}
\vspace{8mm}{\large 21 октября 2026 года\par}
\vfill
\begin{tikzpicture}[x=1cm,y=1cm]
\draw[blue!45!black,line width=0.7pt] (-4,0) .. controls (-2.5,0.6) and (-1.0,-0.6) .. (0,0) .. controls (1.0,0.6) and (2.5,-0.6) .. (4,0);
\end{tikzpicture}\par\vspace{14mm}
\end{titlepage}
\clearpage
\section*{Сводная таблица ошибок}
\small
\renewcommand{\arraystretch}{1.35}
\begin{longtable}{@{}p{0.115\linewidth}p{0.165\linewidth}p{0.27\linewidth}p{0.175\linewidth}p{0.175\linewidth}@{}}
\toprule
№ & Тема & Суть ошибки & Ваш ответ & Правильный ответ\\\midrule\endhead
\hyperref[task:v17]{В1--7} & Логарифмы & Не проверена область допустимых значений & $0{,}5$ & $-0{,}5$\\
\hyperref[task:v19]{В1--9} & Производная и график & Неверно подсчитаны переходы знака производной & 3 & 2\\
\hyperref[task:v114l]{№14 (лог.)} & Тригонометрия & Пропущен корень на отрезке & $-5\pi$ & $-5\pi$, $-11\pi/2$\\
\hyperref[task:v114t]{№14 (триг.)} & Тригонометрия & Неверное разложение на множители & Посторонние значения & Два корня через арктангенс\\
\hyperref[task:v34]{В3--4} & Вероятность & Неверный подсчёт доли подходящих чисел & $0{,}45$ & $0{,}2$\\
\hyperref[task:v35]{В3--5} & Вероятность & Перепутаны события и дополнения & $0{,}6$ & $0{,}14$\\
\hyperref[task:v39]{В3--9} & Производная & Вместо скорости вычислен путь & $14{,}7$ & $3{,}4$ м/с\\
\hyperref[task:v310]{В3--10} & Степени & Перепутано сжатие и расширение & $76{,}8$ л & $0{,}075$ л\\
\hyperref[task:v314]{В3--14} & Показательные уравнения & Неверный дискриминант и отбор & $1;\ 2$ & $1$ на отрезке\\
\hyperref[task:v316]{В3--16} & Логарифм. неравенство & Неправильная смена основания, решение не закончено & Не указан & Объединение трёх промежутков\\
\hyperref[task:v317]{В3--17} & Экстремум функции & Не найдена производная и максимум & Не указан & $\approx510$ м/с\\
\hyperref[task:c2]{Ч--2} & Кредит & 8\% заменены на 80\% & $4500$ тыс. & $1800$ тыс.\\
\hyperref[task:c4]{Ч--4} & Кредит & 15\% приняты за число 15 & $200\%$ & $2\%$ в месяц\\
\hyperref[task:c8]{Ч--8} & Понижение степени & Формула для синуса применена к косинусу & Неверное семейство & $\pm\pi/18+\pi n/3$\\
\hyperref[task:c9]{Ч--9} & Понижение степени & Формула для косинуса применена с неверным знаком & $\pi/8+\pi n/4$ & $\pi/4+\pi n/2$\\
\hyperref[task:c10]{Ч--10} & Тригонометрия & Потерян знак при преобразовании тангенса & Противоречивый переход & $2\pi n$ или $\pi/2+2\pi n$\\
\bottomrule
\end{longtable}
\normalsize
\clearpage
\vspace*{0.38\textheight}\begin{center}\Huge\bfseries Разбор ошибок\end{center}

\Task{v17}{Вариант 1. Задание 7}
\Part{Текст задания}
Найдите больший корень уравнения
\[\log_5(0{,}25-x)=\log_5(x^2-x).\]
\Part{Ваше решение и Ваш ответ}
Вы приравняли аргументы, получили $x^2=0{,}25$ и два значения $x=0{,}5$ и $x=-0{,}5$. В ответ записано $0{,}5$.
\Part{Ошибка}
\Err{не проверена область допустимых значений логарифмов.}
\Part{Где именно возникает ошибка}
До нахождения двух кандидатов ход верен. Первый неверный шаг --- выбор числа $0{,}5$ как допустимого корня без подстановки в исходное уравнение.
\Part{Почему это неверно}
Логарифм существует только от строго положительного числа. При $x=0{,}5$ первый аргумент $0{,}25-0{,}5=-0{,}25$ отрицателен. Второй аргумент $0{,}5^2-0{,}5=-0{,}25$ также отрицателен. Равенство двух отрицательных аргументов не делает логарифмы определёнными. При решении логарифмического уравнения важно одновременно учитывать алгебраическое равенство и область допустимых значений.
\Part{Ключевая идея}
\Idea{сначала выписать условия положительности аргументов, а затем отобрать из найденных кандидатов допустимые.}
\Part{Исправление}
Сохраняем найденные кандидаты, но проверяем каждый из них: $0{,}5$ исключается; при $x=-0{,}5$ оба аргумента равны $0{,}75>0$.
\Part{Полное правильное решение}
\[
\begin{cases}0{,}25-x>0,\\x^2-x>0,\\0{,}25-x=x^2-x.\end{cases}
\]
Из последнего уравнения $x^2=0{,}25$, то есть $x=\pm0{,}5$. Условие $0{,}25-x>0$ означает $x<0{,}25$; условие $x(x-1)>0$ означает $x<0$ или $x>1$. Совместно получаем $x<0$. Следовательно, остаётся только $x=-0{,}5$.
\Part{Проверка}
\[0{,}25-(-0{,}5)=0{,}75,\qquad (-0{,}5)^2-(-0{,}5)=0{,}75.\]
Оба логарифма определены и равны.
\Part{Итог}
\Result{единственный допустимый, а значит и больший корень --- $-0{,}5$.}
\Part{Задача на закрепление}
Найдите больший допустимый корень уравнения $\log_3(1-x)=\log_3(x^2+x)$.

\Task{v19}{Вариант 1. Задание 9}
\Part{Текст задания}
По графику $y=f'(x)$ найдите количество точек максимума функции $f(x)$ на отрезке $[-9;3]$.
\Part{Ваше решение и Ваш ответ}
На графике выделен рассматриваемый отрезок; в ответ записано $3$.
\Part{Ошибка}
\Err{неверно подсчитаны точки, в которых производная меняет знак с плюса на минус.}
\Part{Где именно возникает ошибка}
Вы правильно опираетесь на график производной, но при подсчёте включаете лишнюю точку. На данном отрезке подходящих переходов два.
\Part{Почему это неверно}
Нулевое значение производной само по себе не гарантирует максимум. Если $f'(x)$ меняет знак с отрицательного на положительный, то функция сначала убывает, а потом возрастает: это минимум. Если знак не меняется, экстремума по одному лишь нулю тоже не возникает. Для максимума нужен именно переход $+\to-$.
\Part{Ключевая идея}
\Idea{определять максимумы по направлению смены знака производной, а не по общему числу пересечений графика с осью.}
\Part{Исправление}
На заданном фрагменте графика выделяем только пересечения оси, после которых $f'(x)<0$, а до которых $f'(x)>0$. Их два (приблизительно при $x=-6$ и $x=1$).
\Part{Полное правильное решение}
Просматриваем график слева направо в пределах $[-9;3]$. При первом интересующем пересечении производная меняется с плюса на минус: получаем максимум. Следующее пересечение с переходом минуса в плюс даёт минимум, поэтому в ответ не входит. При следующем переходе плюса в минус появляется второй максимум. Иных переходов $+\to-$ внутри указанного отрезка нет.
\Part{Проверка}
Сопоставляем найденные две точки со знаком производной непосредственно слева и справа от них: оба раза знак меняется $+\to-$.
\Part{Итог}
\Result{на указанном отрезке две точки максимума.}
\Part{Задача на закрепление}
Производная непрерывной дифференцируемой функции меняет знак на прямой в точках $-3$, $0$ и $4$ соответственно по схемам $+\to-$, $-\to+$, $+\to-$. Найдите число точек максимума на отрезке $[-5;5]$.

\Task{v114l}{Задание 14. Логарифмическое уравнение}
\Part{Текст задания}
Решите уравнение
\[\log_2\!\left(\sqrt2\cos\left(x-\frac\pi2\right)-\cos\left(x-\frac\pi4\right)\right)=-\frac12\]
и найдите все его корни на отрезке $[-6\pi;-9\pi/2]$.
\Part{Ваше решение и Ваш ответ}
Вы правильно перешли к уравнению $\sin x-\cos x=1$, затем получили две серии: $x=\pi+2\pi n$ и $x=\pi/2+2\pi n$. При отборе корней на отрезке записано только $-5\pi$.
\Part{Ошибка}
\Err{не выполнен отбор корней второй серии.}
\Part{Где именно возникает ошибка}
Последний правильный шаг --- получение обеих общих серий. Первый неполный шаг --- проверка ограничения на $x$ только для серии $x=\pi+2\pi n$.
\Part{Почему это неверно}
Когда уравнение раскладывается на два множителя, решение представляет объединение корней каждого множителя. Сужение к отрезку надо проводить отдельно для каждой серии. Даже если первая серия уже дала корень, это не исключает корни из второй.
\Part{Ключевая идея}
\Idea{решить оба множителя и независимо подставить в границы отрезка каждую серию корней.}
\Part{Исправление}
Для второй серии проверяем
\[-6\pi\le\frac\pi2+2\pi n\le-\frac{9\pi}{2}.\]
После деления на $\pi$ и преобразований получаем $-13/4\le n\le-5/2$, поэтому $n=-3$ и $x=-11\pi/2$.
\Part{Полное правильное решение}
Так как $2^{-1/2}=\sqrt2/2$, из логарифмического уравнения следует
\[\sqrt2\sin x-\frac{\sqrt2}{2}(\cos x+\sin x)=\frac{\sqrt2}{2}.\]
Умножаем на $\sqrt2$: $\sin x-\cos x=1$. По формулам половинного аргумента:
\[2\sin\frac x2\cos\frac x2-\left(\cos^2\frac x2-\sin^2\frac x2\right)=\sin^2\frac x2+\cos^2\frac x2.\]
Приводим подобные члены:
\[2\cos\frac x2\left(\sin\frac x2-\cos\frac x2\right)=0.\]
Отсюда $x=\pi+2\pi n$ или $x=\pi/2+2\pi n$, где $n\in\mathbb Z$. В первой серии ограничения дают $n=-3$, то есть $-5\pi$; во второй серии также $n=-3$, то есть $-11\pi/2$.
\Part{Проверка}
Для обоих найденных значений $\sin x-\cos x=1$, следовательно, исходный аргумент логарифма равен $\sqrt2/2>0$ и логарифм равен $-1/2$. Оба значения принадлежат данному отрезку.
\Part{Итог}
\Result{на отрезке два корня: $-11\pi/2$ и $-5\pi$.}
\Part{Задача на закрепление}
Решите $\sin x+\cos x=1$ и найдите все корни на отрезке $[-4\pi;-2\pi]$.

\Task{v114t}{Задание 14. Уравнение с двойным углом}
\Part{Текст задания}
Решите $24\sin2x-7\cos2x=15$ и найдите корни на $[-3\pi;-2\pi]$.
\Part{Ваше решение и Ваш ответ}
Вы перешли к записи $48\sin x\cos x-22\cos^2x-8\sin^2x=0$, но затем неверно вынесли множитель, приравняли отдельные части к нулю и получили значения, не следующие из исходного уравнения.
\Part{Ошибка}
\Err{сумма квадратов и смешанного произведения разложена на множители неверно.}
\Part{Где именно возникает ошибка}
До строки $48\sin x\cos x-22\cos^2x-8\sin^2x=0$ преобразования верны. Следующая строка, где всё выражение представлено как произведение с множителем $8\sin x\cos x$, не является тождеством: например, член $-22\cos^2x$ не содержит множителя $\sin x$.
\Part{Почему это неверно}
Выносить общий множитель можно, только если он действительно присутствует в каждом слагаемом. Нельзя применять правило нулевого произведения к сумме, которая не была правильно разложена. В этой задаче удобно переходить к тангенсу, но сначала отдельно проверить случай $\cos x=0$, чтобы деление не потеряло корни.
\Part{Ключевая идея}
\Idea{проверить $\cos x=0$, затем разделить корректное однородное уравнение второй степени на $\cos^2x$.}
\Part{Исправление}
При $\cos x=0$ левая часть исходного уравнения равна $7$, а не $15$; такие $x$ не являются корнями. Поэтому деление на $\cos^2x$ допустимо для всех корней.
\Part{Полное правильное решение}
\[24(2\sin x\cos x)-7(\cos^2x-\sin^2x)=15(\sin^2x+\cos^2x),\]
\[48\sin x\cos x-22\cos^2x-8\sin^2x=0.\]
Делим на $\cos^2x$ и обозначаем $t=\tan x$:
\[48t-22-8t^2=0\quad\Longleftrightarrow\quad 4t^2-24t+11=0.\]
Дискриминант $D=24^2-4\cdot4\cdot11=400$, $\sqrt D=20$, поэтому
\[t=\frac{24-20}{8}=\frac12\quad\text{или}\quad t=\frac{24+20}{8}=\frac{11}{2}.\]
Общие решения:
\[x=\arctan\frac12+\pi n\quad\text{или}\quad x=\arctan\frac{11}{2}+\pi n,\qquad n\in\mathbb Z.\]
Поскольку оба арктангенса лежат между $0$ и $\pi/2$, на отрезке $[-3\pi;-2\pi]$ в обеих сериях подходит только $n=-3$.
\Part{Проверка}
Подстановка $t=1/2$ и $t=11/2$ в квадратное уравнение даёт нуль. При обратной подстановке использованы формулы $\sin2x=2t/(1+t^2)$, $\cos2x=(1-t^2)/(1+t^2)$; они возвращают исходное равенство. Оба выбранных значения расположены внутри требуемого отрезка.
\Part{Итог}
\Result{нужные корни: $-3\pi+\arctan(1/2)$ и $-3\pi+\arctan(11/2)$.}
\Part{Задача на закрепление}
Решите $3\sin2x+4\cos2x=1$ методом перехода к тангенсу, предварительно проверив случай $\cos x=0$.

\Task{v34}{Вариант 3. Задание 4}
\Part{Текст задания}
Из множества натуральных чисел от $12$ до $36$ включительно наудачу выбирают одно число. Найдите вероятность того, что выбранное число делится на $6$.
\Part{Ваше решение и Ваш ответ}
На листе отмечены пять благоприятных исходов и записано $5/25$, однако затем получен неверный десятичный результат $0{,}45$.
\Part{Ошибка}
\Err{неверный перевод обыкновенной дроби в десятичную.}
\Part{Где именно возникает ошибка}
Последний правильный шаг --- дробь $5/25$. Первый неправильный шаг --- равенство $5/25=0{,}45$.
\Part{Почему это неверно}
Числитель показывает число благоприятных исходов, знаменатель --- число всех равновозможных исходов. Значение отношения можно находить сокращением дроби, делением или приведением знаменателя к ста. Здесь $5/25=1/5=20/100$, а не $45/100$. Ошибка возникает уже после правильного подсчёта исходов.
\Part{Ключевая идея}
\Idea{сначала корректно пересчитать все исходы и затем сократить дробь до вычисления десятичной записи.}
\Part{Исправление}
Сохраняем $5/25$ и сокращаем числитель и знаменатель на $5$.
\Part{Полное правильное решение}
Всего целых чисел от $12$ до $36$ включительно $36-12+1=25$. Делятся на $6$ именно $12$, $18$, $24$, $30$ и $36$: таких чисел $5$. Поэтому по классическому определению вероятности
\[P=\frac{5}{25}=\frac15=0{,}2.\]
\Part{Проверка}
Каждый из пяти перечисленных вариантов действительно делится на $6$; других кратных $6$ внутри заданного отрезка нет. Доля $5$ из $25$ равна одной пятой.
\Part{Итог}
\Result{вероятность равна $0{,}2$.}
\Part{Задача на закрепление}
Из всех натуральных чисел от $15$ до $54$ включительно случайно выбирают одно. Найдите вероятность того, что оно делится на $5$.

\Task{v35}{Вариант 3. Задание 5}
\Part{Текст задания}
Вероятность того, что температура окажется выше $+22^\circ\mathrm{C}$, равна $0{,}27$. Вероятность того, что температура окажется ниже $+25^\circ\mathrm{C}$, равна $0{,}87$. Найдите вероятность попадания температуры в промежуток от $+22^\circ\mathrm{C}$ до $+25^\circ\mathrm{C}$.
\Part{Ваше решение и Ваш ответ}
Записано $0{,}87-0{,}27=0{,}6$.
\Part{Ошибка}
\Err{вычтены вероятности двух событий, которые не являются вложенными в требуемом направлении.}
\Part{Где именно возникает ошибка}
Первый неверный переход --- принятие события «температура выше $22$» за событие, которое нужно вычесть из «температура ниже $25$». На самом деле вычитать нужно событие «температура не выше $22$».
\Part{Почему это неверно}
События «выше $22$» и «ниже $25$» сильно пересекаются: их пересечение и есть искомая область. При вычитании $0{,}87-0{,}27$ складывается неверное разбиение. Нужно сначала дополнить событие «выше $22$» до единицы, получив вероятность нахождения слева от $22$.
\Part{Ключевая идея}
\Idea{представить искомую область как «ниже 25» за вычетом «не выше 22».}
\Part{Исправление}
\[P(T\le22)=1-P(T>22)=1-0{,}27=0{,}73.\]
\Part{Полное правильное решение}
Для температуры как непрерывной величины вероятность точного значения границы равна нулю. Поэтому
\[\begin{aligned}
P(22<T<25)&=P(T<25)-P(T\le22)\\
&=0{,}87-(1-0{,}27)\\
&=0{,}87-0{,}73=0{,}14.
\end{aligned}\]
\Part{Проверка}
Область ниже $25$ разбивается на две непересекающиеся части: температура не выше $22$ с вероятностью $0{,}73$ и промежуток между границами с вероятностью $0{,}14$. Их сумма $0{,}87$, как в условии.
\Part{Итог}
\Result{искомая вероятность равна $0{,}14$.}
\Part{Задача на закрепление}
Вероятность температуры выше $18^\circ\mathrm{C}$ равна $0{,}4$, а ниже $24^\circ\mathrm{C}$ --- $0{,}85$. Найдите вероятность температуры между этими значениями.

\Task{v39}{Вариант 3. Задание 9}
\Part{Текст задания}
Материальная точка движется по закону $x(t)=0{,}2t^2+t+1{,}5$, где расстояние измеряется в метрах, время --- в секундах. Найдите скорость в момент $t=6$.
\Part{Ваше решение и Ваш ответ}
Подставлено $t=6$ в закон координаты и получено $x(6)=14{,}7$; число $14{,}7$ записано как ответ.
\Part{Ошибка}
\Err{вместо мгновенной скорости вычислена координата точки.}
\Part{Где именно возникает ошибка}
Подстановка в $x(t)$ и арифметический расчёт координаты выполнены правильно. Ошибка начинается при отождествлении значения $x(6)$ со скоростью $v(6)$.
\Part{Почему это неверно}
Координата показывает положение точки, а скорость --- насколько быстро изменяется координата в данный момент. Это разные величины, даже размерности отличаются: метры и метры в секунду. При заданном законе движения мгновенная скорость равна производной координаты по времени, то есть $v(t)=x'(t)$.
\Part{Ключевая идея}
\Idea{сначала найти производную закона движения, затем подставить указанный момент времени.}
\Part{Исправление}
Сохраняем заданную формулу $x(t)$, но вместо вычисления её значения находим $x'(t)$.
\Part{Полное правильное решение}
Правило производной степени даёт $(0{,}2t^2)'=0{,}4t$; производная $t$ равна $1$, а постоянной $1{,}5$ --- нулю. Значит,
\[v(t)=x'(t)=0{,}4t+1,\qquad v(6)=0{,}4\cdot6+1=3{,}4\ \text{м/с}.\]
\Part{Проверка}
На секунду позже скорость увеличивается на $0{,}4$ м/с, что согласуется с положительным ускорением $x''(t)=0{,}4$. Размерность ответа --- м/с.
\Part{Итог}
\Result{мгновенная скорость через шесть секунд равна $3{,}4$ м/с.}
\Part{Задача на закрепление}
Точка движется по закону $x(t)=0{,}3t^2+2t-4$. Найдите её мгновенную скорость в момент $t=5$.

\Task{v310}{Вариант 3. Задание 10}
\Part{Текст задания}
При адиабатическом сжатии газа выполняется закон $p_1V_1^{1{,}4}=p_2V_2^{1{,}4}$. Начальные параметры: $p_1=0{,}5$ атмосферы, $V_1=2{,}4$ литра. Найдите конечный объём при $p_2=64$ атмосферы.
\Part{Ваше решение и Ваш ответ}
Вы верно записали адиабатический закон, но в дальнейшем вместо деления начального объёма на $32$ умножили его на $32$ и получили $76{,}8$ литра.
\Part{Ошибка}
\Err{перевёрнуто отношение давлений при извлечении степени.}
\Part{Где именно возникает ошибка}
Последний верный шаг --- равенство $0{,}5\cdot2{,}4^{1{,}4}=64\cdot V_2^{1{,}4}$. Неверный шаг --- преобразование, приводящее к $V_2=2{,}4\cdot32$, хотя давление выросло в $128$ раз.
\Part{Почему это неверно}
Если произведение $pV^{1{,}4}$ постоянно, то при увеличении давления объём должен уменьшиться. При выражении нового объёма необходимо брать отношение начального давления к конечному: $p_1/p_2$, а не обратное. Проверка физического смысла уже показывает, что ответ $76{,}8$ л невозможен при сжатии.
\Part{Ключевая идея}
\Idea{выразить новый объём через отношение $p_1/p_2$ и проследить направление изменения величин.}
\Part{Исправление}
\[V_2=V_1\left(\frac{p_1}{p_2}\right)^{1/1{,}4}=2{,}4\left(\frac1{128}\right)^{5/7}.\]
\Part{Полное правильное решение}
Поскольку $1{,}4=7/5$, обратная степень равна $5/7$. Кроме того, $128=2^7$, поэтому
\[\left(\frac1{128}\right)^{5/7}=(2^{-7})^{5/7}=2^{-5}=\frac1{32}.\]
Следовательно, $V_2=2{,}4/32=0{,}075$ литра.
\Part{Проверка}
\[\frac{p_2}{p_1}=128,\qquad\left(\frac{V_1}{V_2}\right)^{1{,}4}=32^{7/5}=(2^5)^{7/5}=128.\]
Обе стороны адиабатического закона согласованы. Новый объём меньше начального.
\Part{Итог}
\Result{конечный объём равен $0{,}075$ литра.}
\Part{Задача на закрепление}
Для газа выполняется $p_1V_1^{1{,}5}=p_2V_2^{1{,}5}$. При $p_1=1$, $V_1=8$ л и $p_2=8$ найдите $V_2$.

\Task{v314}{Вариант 3. Задание 14}
\Part{Текст задания}
Решите уравнение
\[4^{x^3-x^2+3x-0{,}5}+4^{x^3+x^2-3x+2{,}5}-5\cdot2^{2x^3+1}=0\]
и найдите его корни на отрезке $[0{,}4;1{,}4]$.
\Part{Ваше решение и Ваш ответ}
Вы выполнили показательную замену $t=4^{x^2-3x}$ и верно получили $t=1/4$ или $t=1/16$. Однако из уравнения $x^2-3x+1=0$ ошибочно получены корни $1$ и $2$; итог записан как $2;1$ без надлежащего отбора по отрезку.
\Part{Ошибка}
\Err{неверный дискриминант квадратного уравнения и неполный отбор корней по условию.}
\Part{Где именно возникает ошибка}
Правильный последний этап --- два значения $t$. Первый неверный шаг --- вычисление $D=1$ для уравнения $x^2-3x+1=0$. В действительности $D=9-4=5$.
\Part{Почему это неверно}
Разные значения замены порождают разные квадратные уравнения. Для первого нужно $x^2-3x=-1$, а для второго $x^2-3x=-2$. Их решения нельзя смешивать. После получения всех корней обязательно нужно проверить заданный отрезок: число $2$ в него не входит.
\Part{Ключевая идея}
\Idea{разделять две ветви обратной замены и поочерёдно проверять каждое найденное число.}
\Part{Исправление}
Для $t=1/4$ получаем $x=(3\pm\sqrt5)/2$; для $t=1/16$ --- $x=1$ или $x=2$.
\Part{Полное правильное решение}
Делим исходное уравнение на положительное число $2^{2x^3+1}$. После приведения степеней получаем
\[\frac1{4\cdot4^{x^2-3x}}+16\cdot4^{x^2-3x}-5=0.\]
Положим $t=4^{x^2-3x}>0$. Тогда
\[\frac1{4t}+16t-5=0\iff64t^2-20t+1=0.\]
Дискриминант равен $400-256=144$, поэтому
\[t=\frac{20\pm12}{128}\in\left\{\frac14,\frac1{16}\right\}.\]
При $t=1/4$ имеем $x^2-3x=-1$, откуда
\[x=\frac{3\pm\sqrt5}{2}.\]
При $t=1/16$ имеем $x^2-3x=-2$, откуда $x=1$ или $x=2$. Приближённо $(3-\sqrt5)/2\approx0{,}382<0{,}4$, а $(3+\sqrt5)/2\approx2{,}618>1{,}4$; $2>1{,}4$. На указанном отрезке остаётся только $1$.
\Part{Проверка}
При $x=1$ замена даёт $4^{1-3}=1/16$; подстановка в преобразованное уравнение даёт $1/(4/16)+16/16-5=4+1-5=0$. Точка $x=1$ принадлежит $[0{,}4;1{,}4]$.
\Part{Итог}
\Result{все действительные корни: $1$, $2$, $(3-\sqrt5)/2$, $(3+\sqrt5)/2$; на заданном отрезке только $1$.}
\Part{Задача на закрепление}
Решите $(2^{x^2-4x})^2-5\cdot2^{x^2-4x}+4=0$ и выберите корни из $[0;3]$.

\Task{v316}{Вариант 3. Задание 16}
\Part{Текст задания}
Решите неравенство
\[\log_{\sqrt2}^{4}(24-x^2)-192\log_4^{2}(24-x^2)\ge64.\]
\Part{Ваше решение и Ваш ответ}
Вы записали условие $24-x^2>0$ и предприняли переход к логарифмам с основанием $2$, после чего получили квадратное неравенство с коэффициентом $-48$ при замене, но не завершили обратную подстановку и не указали ответ.
\Part{Ошибка}
\Err{при приведении логарифмов к одному основанию потерян множитель $16$ и не завершена работа с заменой.}
\Part{Где именно возникает ошибка}
Верный шаг --- ОДЗ $-\sqrt{24}<x<\sqrt{24}$. Неверный --- представление первой четвёртой степени логарифма без учёта того, что $\log_{\sqrt2}u=2\log_2u$ и потому четвёртая степень увеличивает выражение в $2^4=16$ раз.
\Part{Почему это неверно}
Показатель степени относится ко всему значению логарифма. Коэффициент перехода между основаниями возводится в ту же степень. Для второго логарифма основание $4$ даёт множитель $1/2$, но после возведения в квадрат этот множитель становится $1/4$. Без обоих коэффициентов квадратное неравенство будет другим.
\Part{Ключевая идея}
\Idea{последовательно применить формулу смены основания, затем решить неравенство относительно квадрата логарифма и возвратиться к $x$.}
\Part{Исправление}
При $u=24-x^2>0$ положим $y=\log_2u$. Тогда
\[\log_{\sqrt2}u=2y,\qquad\log_4u=\frac y2.\]
Следовательно, исходное неравенство равно $16y^4-48y^2\ge64$.
\Part{Полное правильное решение}
\[16y^4-48y^2-64\ge0\iff y^4-3y^2-4\ge0\iff(y^2-4)(y^2+1)\ge0.\]
Поскольку $y^2+1>0$, остаётся $y^2\ge4$, то есть $y\ge2$ либо $y\le-2$.

Если $y\ge2$, то $24-x^2\ge4$, поэтому $x^2\le20$, или $-2\sqrt5\le x\le2\sqrt5$.

Если $y\le-2$, то $0<24-x^2\le1/4$. Из этого следует $95/4\le x^2<24$, то есть
\[x\in(-2\sqrt6;-\sqrt{95}/2]\ \cup\ [\sqrt{95}/2;2\sqrt6).\]
Объединяем обе ветви.
\Part{Проверка}
На границах $x=\pm2\sqrt5$ аргумент равен $4$ и $y=2$; на границах $x=\pm\sqrt{95}/2$ аргумент равен $1/4$ и $y=-2$. При $x=\pm2\sqrt6$ аргумент равен нулю и логарифм не существует, поэтому концы не включены.
\Part{Итог}
\Result{$x\in(-2\sqrt6;-\sqrt{95}/2]\cup[-2\sqrt5;2\sqrt5]\cup[\sqrt{95}/2;2\sqrt6)$.}
\Part{Задача на закрепление}
Решите $\log_{\sqrt2}^{4}(9-x^2)-48\log_4^{2}(9-x^2)\ge16$.

\Task{v317}{Вариант 3. Задание 17}
\Part{Текст задания}
Распределение скорости молекул кислорода описывается функцией Максвелла
\[f(v)=4\pi v^2\!\left(\frac M{2\pi RT}\right)^{3/2}e^{-Mv^2/(2RT)},\quad v\ge0.\]
Найдите скорость максимума плотности при $T=500$ К, $M=0{,}032$ кг/моль и $R\approx8{,}31$ Дж/(моль$\cdot$К).
\Part{Ваше решение и Ваш ответ}
В решении предпринята попытка дифференцировать экспоненциальный множитель, но корректная производная произведения не получена; поиск критических точек и ответ отсутствуют.
\Part{Ошибка}
\Err{не использовано правило производной произведения и цепное правило для показательной функции.}
\Part{Где именно возникает ошибка}
Правильно выписана сама функция $f(v)$. Первый неполный шаг --- запись $f'(v)$ без производной одновременно множителя $v^2$ и показателя $-Mv^2/(2RT)$.
\Part{Почему это неверно}
В функции перемножаются степенной множитель и экспонента, причём оба зависят от скорости. Производная произведения равна сумме двух слагаемых. У экспоненты производная равна самой экспоненте, умноженной на производную её показателя. Если отбросить хотя бы один вклад, критическая точка получится неверной либо не будет найдена.
\Part{Ключевая идея}
\Idea{вынести положительные постоянные, продифференцировать $v^2e^{-av^2}$ и исследовать знак производной.}
\Part{Исправление}
Обозначим $C=4\pi(M/(2\pi RT))^{3/2}>0$ и $a=M/(2RT)>0$. Тогда $f(v)=Cv^2e^{-av^2}$.
\Part{Полное правильное решение}
Используя правило произведения,
\[\begin{aligned}
f'(v)&=C\bigl(2ve^{-av^2}+v^2e^{-av^2}(-2av)\bigr)\\
&=2Cve^{-av^2}(1-av^2).
\end{aligned}\]
При $v>0$ первые три множителя положительны, поэтому знак $f'(v)$ определяется $1-av^2$. Производная положительна при $v<1/\sqrt a$, равна нулю при $v=1/\sqrt a$ и отрицательна при $v>1/\sqrt a$. Следовательно, это точка максимума. На краю области $f(0)=0$, а для положительных скоростей функция положительна.
\[v_{\max}=\frac1{\sqrt a}=\sqrt{\frac{2RT}{M}}=\sqrt{\frac{2\cdot8{,}31\cdot500}{0{,}032}}\approx509{,}6\ \text{м/с}.\]
\Part{Проверка}
Слева от найденной скорости функция возрастает, справа --- убывает. Проверка размерности подкоренного выражения даёт $\mathrm{м}^2/\mathrm{с}^2$, значит результат измеряется в м/с.
\Part{Итог}
\Result{скорость, при которой плотность максимальна, составляет приблизительно $510$ м/с.}
\Part{Задача на закрепление}
Для $v\ge0$ найдите точку максимума функции $g(v)=v^2e^{-v^2/18}$.

\Task{c2}{Блок «Чек». Задание 2}
\Part{Текст задания}
Кредит $1500$ тыс. руб. взят на четыре года под $8\%$ годовых по дифференцированной схеме. Найдите общую сумму выплат и переплату.
\Part{Ваше решение и Ваш ответ}
Вы верно выбрали формулу переплаты $S\,p(N+1)/2$, но при вычислении подставили $0{,}8$ вместо $0{,}08$. Записаны переплата $3000$ тыс. руб. и общая сумма $4500$ тыс. руб.
\Part{Ошибка}
\Err{неправильный перевод восьми процентов в десятичную дробь.}
\Part{Где именно возникает ошибка}
Последний правильный шаг --- формула $S\,p(N+1)/2$. Неверный переход --- использование $p=0{,}8$, что означает восемьдесят процентов, а не восемь.
\Part{Почему это неверно}
Процент --- сотая доля целого: $8\%=8/100=0{,}08$. Число $0{,}8=80/100$ соответствует $80\%$ и даёт переплату в десять раз больше. При равномерном уменьшении основного долга ежегодные проценты начисляются на остатки $S$, $3S/4$, $S/2$ и $S/4$, поэтому сумма процентов равна $S\,p(1+2+3+4)/4$.
\Part{Ключевая идея}
\Idea{сразу записать процентную ставку как $p=r/100$ и учитывать остаток долга в каждом году.}
\Part{Исправление}
Сохраняем используемую формулу и подставляем $p=0{,}08$, а не $0{,}8$.
\Part{Полное правильное решение}
Ежегодно гасится основная сумма $1500/4=375$ тыс. руб. Проценты по годам:
\[1500\cdot0{,}08=120;\quad1125\cdot0{,}08=90;\quad750\cdot0{,}08=60;\quad375\cdot0{,}08=30.\]
Переплата составляет $120+90+60+30=300$ тыс. руб. Общая сумма выплат --- основной долг плюс проценты:
\[1500+300=1800\ \text{тыс. руб.}\]
Тот же результат даёт формула $1500\cdot0{,}08\cdot5/2=300$.
\Part{Проверка}
Сумма четырёх долей основного долга равна $4\cdot375=1500$, сумма процентов --- $300$, итог --- $1800$.
\Part{Итог}
\Result{переплата $300$ тыс. руб.; все выплаты $1800$ тыс. руб.}
\Part{Задача на закрепление}
Кредит $2400$ тыс. руб. взят на три года под $6\%$ годовых с равномерным погашением основного долга. Найдите общую сумму выплат и переплату.

\Task{c4}{Блок «Чек». Задание 4}
\Part{Текст задания}
Кредит с равномерным ежемесячным погашением основного долга взят на $14$ месяцев. Общая сумма платежей больше суммы кредита на $15\%$. Найдите месячную процентную ставку.
\Part{Ваше решение и Ваш ответ}
Составлено соотношение для переплаты с множителем $(14+1)/2$, но в правой части записано $15$ вместо $0{,}15$. Отсюда найдено $p=2$ и дан ответ $200\%$.
\Part{Ошибка}
\Err{процентная величина подставлена как число, а не как доля.}
\Part{Где именно возникает ошибка}
Правильный шаг --- понимание, что общая переплата равна сумме процентов за все месяцы. Первый неверный шаг --- равенство $15=15p/2$; сумма переплаты на самом деле составляет $0{,}15S$.
\Part{Почему это неверно}
Фраза «превышает на пятнадцать процентов» означает добавку $15/100$ от суммы кредита. Формула процентов для равномерного погашения основного долга пропорциональна самому кредиту $S$, поэтому после сокращения $S$ остаётся доля $0{,}15$, а не процентное число $15$.
\Part{Ключевая идея}
\Idea{записывать общую переплату как $0{,}15S$ и только потом сокращать сумму кредита.}
\Part{Исправление}
\[0{,}15=\frac{15p}{2},\qquad p=\frac{0{,}30}{15}=0{,}02.\]
\Part{Полное правильное решение}
При сроке $N=14$ месяцев остатки основного долга перед начислением процентов равны $S$, $13S/14$, $12S/14$, \ldots, $S/14$. Поэтому общая сумма процентов равна
\[Sp\left(1+\frac{13}{14}+\cdots+\frac1{14}\right)=Sp\frac{1+14}{2}=\frac{15Sp}{2}.\]
По условию эта сумма равна $0{,}15S$. Делим обе стороны на $S>0$:
\[\frac{15p}{2}=0{,}15\quad\Longrightarrow\quad p=0{,}02=2\%.\]
\Part{Проверка}
Если ежемесячная ставка $2\%$, то переплата составляет $S\cdot0{,}02\cdot15/2=0{,}15S$, то есть ровно $15\%$ первоначального кредита.
\Part{Итог}
\Result{месячная ставка равна $2\%$.}
\Part{Задача на закрепление}
Кредит с равномерным уменьшением основного долга выдан на девять месяцев. Все проценты составили $10\%$ суммы кредита. Найдите месячную ставку.

\Task{c8}{Блок «Чек». Задание 8}
\Part{Текст задания}
Решите уравнение $\cos^2(3x)=\dfrac34$ с использованием формулы понижения степени.
\Part{Ваше решение и Ваш ответ}
В решении записана формула вида $(1-\cos6x)/2$, применимая к квадрату синуса, и вместо заданной дроби $3/4$ в одном из переходов появляется $3/5$. Далее получено семейство корней, не удовлетворяющее исходному уравнению.
\Part{Ошибка}
\Err{перепутаны формулы понижения степени для синуса и косинуса; дополнительно искажена правая часть.}
\Part{Где именно возникает ошибка}
Первый неверный шаг --- запись $\cos^2(3x)=(1-\cos6x)/2$. Верная формула имеет знак плюс.
\Part{Почему это неверно}
Из тождества $\cos2a=2\cos^2a-1$ следует $\cos^2a=(1+\cos2a)/2$. Тождество $\sin^2a=(1-\cos2a)/2$ отличается знаком. Подстановка одной формулы вместо другой превращает исходное уравнение в уравнение для другой функции. Кроме того, число $3/4$ из условия нельзя менять на $3/5$.
\Part{Ключевая идея}
\Idea{отдельно помнить два тождества понижения степени и внимательно переносить правую часть.}
\Part{Исправление}
\[\frac{1+\cos6x}{2}=\frac34\quad\Longrightarrow\quad1+\cos6x=\frac32\quad\Longrightarrow\quad\cos6x=\frac12.\]
\Part{Полное правильное решение}
Основные решения $\cos u=1/2$ имеют вид $u=\pm\pi/3+2\pi n$ при $n\in\mathbb Z$. Полагая $u=6x$, получаем
\[6x=\pm\frac\pi3+2\pi n\quad\Longrightarrow\quad x=\pm\frac\pi{18}+\frac{\pi n}{3},\qquad n\in\mathbb Z.\]
Знак $\pm$ представляет две серии, которые вместе охватывают все корни.
\Part{Проверка}
При найденных $x$ значение $\cos6x=1/2$, поэтому $\cos^2(3x)=(1+1/2)/2=3/4$, как требуется.
\Part{Итог}
\Result{$x=\pm\pi/18+\pi n/3$, $n\in\mathbb Z$.}
\Part{Задача на закрепление}
Решите уравнение $\cos^2(4x)=\dfrac14$, используя формулу понижения степени.

\Task{c9}{Блок «Чек». Задание 9}
\Part{Текст задания}
Решите уравнение
\[4^{2+\cos^2x}-4^{\cos^2x}=9^{1+\cos^2x}+9^{\cos^2x}.\]
\Part{Ваше решение и Ваш ответ}
После верной показательной замены получено $\cos^2x=0{,}5$. Затем использована формула с минусом для квадрата косинуса и записано семейство $x=\pi/8+\pi n/4$.
\Part{Ошибка}
\Err{в тригонометрической части использована неверная формула и неправильно разделён период решения.}
\Part{Где именно возникает ошибка}
Последний правильный шаг --- $\cos^2x=1/2$. Первый неверный шаг --- подмена $\cos^2x=(1-\cos2x)/2$; правильный знак плюс. Дополнительно из $\cos2x=0$ следует $x=\pi/4+\pi n/2$, а не $\pi/8+\pi n/4$.
\Part{Почему это неверно}
Формула с минусом относится к $\sin^2x$. В данном частном случае она также приводит к $\cos2x=0$, потому что справа стоит $1/2$, но это случайное совпадение, не оправдывающее неверного правила. Независимо от этого ошибка при делении аргумента $2x$ вдвое добавляет посторонние значения.
\Part{Ключевая идея}
\Idea{сначала точно решить показательное уравнение, затем применить верное тождество и аккуратно разделить весь аргумент.}
\Part{Исправление}
Сохраняем $\cos^2x=1/2$ и используем $(1+\cos2x)/2=1/2$. Отсюда $\cos2x=0$.
\Part{Полное правильное решение}
Обозначим $t=\cos^2x\in[0;1]$. Переносим степенные множители:
\[4^t(4^2-1)=9^t(9+1)\quad\Longrightarrow\quad15\cdot4^t=10\cdot9^t.\]
Делим на $15\cdot9^t>0$:
\[\left(\frac49\right)^t=\frac23\quad\Longrightarrow\quad\left(\frac23\right)^{2t}=\left(\frac23\right)^1.\]
Значит, $2t=1$ и $t=1/2$. Тогда $\cos2x=0$, поэтому
\[2x=\frac\pi2+\pi n\quad\Longrightarrow\quad x=\frac\pi4+\frac{\pi n}{2},\qquad n\in\mathbb Z.\]
\Part{Проверка}
В каждом найденном корне $\cos^2x=1/2$; исходное равенство превращается в $15\sqrt4=10\sqrt9$, то есть $30=30$.
\Part{Итог}
\Result{$x=\pi/4+\pi n/2$, $n\in\mathbb Z$.}
\Part{Задача на закрепление}
Решите $2^{2+\cos^2x}-2^{\cos^2x}=3^{1+\cos^2x}+3^{\cos^2x}$.

\Task{c10}{Блок «Чек». Задание 10}
\Part{Текст задания}
Решите уравнение $\sin x+\cos x=1$. Объясните, почему при делении на $\cos x$ необходимо отдельно рассмотреть случай $\cos x=0$.
\Part{Ваше решение и Ваш ответ}
Вы применили формулы половинного угла и правильно получили произведение $2\sin(x/2)(\cos(x/2)-\sin(x/2))=0$. Во второй ветви записано $1-\tan(x/2)=0$, но затем ошибочно сделан вывод $\tan(x/2)=-1$, хотя последующее семейство $x=\pi/2+2\pi n$ соответствует тангенсу $+1$.
\Part{Ошибка}
\Err{ошибка знака при переносе слагаемого; записанный вывод не следует из предыдущего равенства.}
\Part{Где именно возникает ошибка}
Последняя верная строка --- $1-\tan(x/2)=0$. Следующая строка должна быть $\tan(x/2)=1$, а не $-1$.
\Part{Почему это неверно}
При решении равенства $1-t=0$ прибавляем $t$ к обеим частям и получаем $t=1$. Если подставить $t=-1$ назад, левая часть будет $1-(-1)=2$, а не ноль. В итоговом ответе по случайности указано верное семейство, однако математическое обоснование для него нарушено.
\Part{Ключевая идея}
\Idea{сохранять знак при каждом переносе и проверять обе ветви нулевого произведения прямой подстановкой.}
\Part{Исправление}
Из $\cos(x/2)-\sin(x/2)=0$ следует $\tan(x/2)=1$, поскольку при этом равенстве $\cos(x/2)\ne0$.
\Part{Полное правильное решение}
После применения формул двойного угла и переноса единицы:
\[2\sin\frac x2\cos\frac x2+\cos^2\frac x2-\sin^2\frac x2=\sin^2\frac x2+\cos^2\frac x2,\]
\[2\sin\frac x2\left(\cos\frac x2-\sin\frac x2\right)=0.\]
Первая ветвь $\sin(x/2)=0$ даёт $x=2\pi n$. Вторая ветвь $\cos(x/2)=\sin(x/2)$ даёт $x/2=\pi/4+\pi n$, то есть $x=\pi/2+2\pi n$.

Если решать делением исходного уравнения на $\cos x$, случай $\cos x=0$ нужно проверять отдельно: при $x=\pi/2+\pi n$ исходное равенство верно только для чётных $n$. Без этой проверки можно потерять целую серию корней.
\Part{Проверка}
При $x=2\pi n$ получаем $0+1=1$; при $x=\pi/2+2\pi n$ получаем $1+0=1$.
\Part{Итог}
\Result{$x=2\pi n$ или $x=\pi/2+2\pi n$, $n\in\mathbb Z$.}
\Part{Задача на закрепление}
Решите $\sin x-\cos x=1$, используя формулы половинного угла и проверяя условия деления.

\vfill
\begin{center}\small Лёвин Артём Александрович эксклюзивно для Николь Саркисьянц\end{center}
\end{document}